\section{Project Background}
The Norwegian Armed Forces Cyber Defense unit for Computer Network Defense
(CND), Centre for Protection of Critical Infrastructure (Avdeling BKI in
Norwegian), is responsible for detecting and stopping cyber-attacks against the
critical infrastructure and command and control systems utilized by the Norwegian military.
One of the tools used to carry out this task is the open source IDS \snort{}
\cite{Snort}.

\snort{} is a leading \cite{snort_community} open source IDS that is used to
detect potentially unwanted or malicious network traffic based on predefined patterns described in
signatures. There are both free and commercial signatures available for
\snort{} and as \snort{} has been highly adopted by the security community,
these signatures are continually updated and maintained. For large networks, an
IDS will often operate in multiple parts of the network simultaneously, where each separate instance of the IDS is known as a
sensor. 

The management of \snort{} signatures (henceforth referred to as
rules) is cumbersome in multi-sensor environments without a centralized service
which can assure that all sensors are synchronised with the latest rules.
Avdeling BKI has therefore challenged a group of students from Gjøvik University
College (GUC) to design and develop a functional, effective and secure
rule management system for \snort{} that can continue to live on beyond this project.


% Using \snort{} has been plagued for years by the fact that it does not have a good
% system for handling, updating or maintaining the signatures (henceforth called
% rules and rule sets) used for detection, thus making it cumbersome to utilize
% \snort{} for security analysts (ref?).
% 
